% --- Slide 9 (từ part2.tex: Compact Box) ---
\begin{frame}{Compact Box -- đòn bẩy tăng tốc \#1 của FastGS}
\begin{columns}[T]
\begin{column}{0.52\textwidth}
\footnotesize
    \begin{itemize}
        \item Mặc định: bounding-box $3\sigma$ -- khá rộng, chạm nhiều tile không cần thiết
        \item FastGS thay bằng \textbf{compact box}: vùng ảnh hưởng 2D nhỏ hơn, điều chỉnh bằng hệ số \texttt{-{}-mult}
        \item Kết quả: giảm đáng kể số tile mỗi Gaussian phải rasterize
    \end{itemize}
\end{column}
\begin{column}{0.48\textwidth}
    \centering
    \includegraphics[width=0.95\linewidth]{ch3_compact_box.png}
    \captionof{figure}{\footnotesize So sánh bounding-box mặc định và compact box}
\end{column}
\end{columns}
\end{frame}

% --- Slide 10 (từ part2.tex: Alpha blending front-to-back) ---
\begin{frame}{Alpha blending (front-to-back)}
    \begin{itemize}
        \item Màu pixel là tổng có trọng số của các Gaussian chồng nhau theo thứ tự độ sâu
        \item Công thức blend:
        \[
            C = \sum_i c_i\,\alpha_i\,T_i, \qquad T_i = \prod_{j<i}(1-\alpha_j)
        \]
        \item $T_i$: độ truyền sáng còn lại (transmittance) trước Gaussian thứ $i$
    \end{itemize}
    \vspace{0.2cm}
    \centering
    \includegraphics[width=0.55\linewidth]{ch4_alpha_blend_stack.png}
    \captionof{figure}{Chồng lớp alpha blending theo độ sâu}
\end{frame}

% --- Slide 11 (từ part2.tex: Tổng kết phương trình render) ---
\begin{frame}{Tổng kết phương trình render}
    \begin{itemize}
        \item Phương trình render đầy đủ:
        \[
            C(p) = \sum_i c_i\,\alpha_i \prod_{j<i}(1-\alpha_j)
        \]
        \item Khả vi theo mọi tham số: vị trí $\mu$, hiệp phương sai $\Sigma$, độ mờ $\alpha$, hệ số SH
    \end{itemize}
    \vspace{0.4cm}
    \centering
    \begin{tikzpicture}[node distance=1.1cm, every node/.style={font=\footnotesize}]
        \node (g3d) {Gaussian 3D};
        \node (proj) [right=of g3d] {Chiếu 2D};
        \node (alpha) [right=of proj] {Trường alpha};
        \node (blend) [right=of alpha] {Blend};
        \node (pixel) [right=of blend] {Pixel};
        \draw[-Latex] (g3d) -- (proj);
        \draw[-Latex] (proj) -- (alpha);
        \draw[-Latex] (alpha) -- (blend);
        \draw[-Latex] (blend) -- (pixel);
    \end{tikzpicture}
\end{frame}

% --- Slide 12 (từ part3.tex: Hàm mất mát L1 + D-SSIM) ---
\begin{frame}{Hàm mất mát: kết hợp L1 và D-SSIM}
    \begin{columns}[T]
        \begin{column}{0.5\textwidth}
            \begin{itemize}
                \item Loss huấn luyện kết hợp hai thành phần bổ trợ nhau:
                \begin{itemize}
                    \item \textbf{L1}: sai khác trung bình theo từng pixel.
                    \item \textbf{D-SSIM}: $1-\text{SSIM}$, phản ánh sai lệch \textbf{cấu trúc}.
                \end{itemize}
                \item Công thức tổng:
                \[
                    \mathcal{L} = (1-\lambda)\,\mathcal{L}_1 + \lambda\,(1-\text{SSIM})
                \]
                \item Hệ số $\lambda = $ \texttt{lambda\_dssim}: mặc định $0.2$, \texttt{fastgs-lite}: $0.25$
            \end{itemize}
        \end{column}
        \begin{column}{0.5\textwidth}
            \centering
            \includegraphics[width=0.95\linewidth]{ch5_loss_mix.png}
            \captionof{figure}{Tỉ trọng L1 và D-SSIM trong loss}
        \end{column}
    \end{columns}
\end{frame}
